\section{Background}
\label{sec:background}

The Bailout Protocol emerges from three intersecting lineages of research --- accountability in multi-agent AI systems, human-in-the-loop design, and fail-safe design in safety-critical systems --- and is grounded architecturally in the BX3 three-layer model. This section traces each lineage, establishes the specific gaps that motivate the protocol, and defines the three-layer model against which the protocol operates.

% ---------------------------------------------------------------
\subsection{Accountability in Multi-Agent AI Systems}
\label{sec:accountability}

Multi-agent AI systems introduce accountability challenges that do not arise in single-agent deployments. When multiple autonomous or semi-autonomous agents coordinate to decompose and execute complex directives, responsibility for outcomes --- correct, incorrect, or harmful --- becomes distributed across the agent network. This distribution is compounded when agents spawn sub-agents, delegate subtasks laterally, or operate asynchronously under partial observability. The result is an accountability gap: the chain of causal responsibility between an outcome and the original human principal is mediated by multiple machine actors in ways that are opaque, non-linear, and legally ambiguous~\cite{bansal2019}.

Traditional accountability frameworks borrowed from organizational theory presuppose a fixed hierarchy of human agents, each with delineated authority and formal reporting relationships. In software systems, these take the form of access-control lists, role-based security models, and audit logs. However, these frameworks were designed for deterministic, fully-specified systems in which all possible states are enumerated at design time. Modern AI agents --- particularly those operating with large language model backends --- introduce states that are not exhaustively enumerable, whose behaviors emerge from learned statistical patterns rather than explicit rule sets, and whose failure modes are neither fully predictable nor fully auditable post hoc~\cite{dijkstra1982}. A machine actor in such a system may produce an outcome that is neither clearly authorized nor clearly unauthorized within its provisioned scope --- a grey zone that conventional binary permission models cannot capture.

The BX3 Framework addresses this through a structured accountability architecture in which every node carries an immutable parent pointer and a designated human anchor. The parent pointer defines the chain of delegation; the human anchor closes the loop by ensuring that some named human principal is architecturally reachable from every node, regardless of the depth or complexity of the agent network. This is not merely a governance convention --- it is an architectural compulsion: the BX3 state machine does not permit any machine actor to serve as a terminal resolution point for an unresolved exception. The accountability chain is append-only and cryptographically sealed, ensuring that any audit can reconstruct the complete delegation and exception history for any node at any point in time~\cite{wiener1948}.

The relationship between accountability and auditability is central here. Accountability without auditability is a statement of intent; auditability without accountability is a historical record with no enforcement pathway. The BX3 Forensic Ledger provides the latter; the Bailout Protocol, operating through the Purpose and Bounds layers, provides the former. Together they constitute a closed-loop accountability system in which every decision is auditable and every unresolved decision reaches a human~\cite{bansal2019}.

% ---------------------------------------------------------------
\subsection{Human-in-the-Loop Design}
\label{sec:hitl}

Human-in-the-loop (HITL) design positions human agents as integral components of AI-driven workflows, rather than as external reviewers who intervene only after the system has reached a conclusion. The canonical motivation for HITL is epistemic: AI systems, regardless of their statistical sophistication, operate within the bounds of their training data and objective functions and cannot reliably reason about novel situations, value trade-offs, or contextual factors that require human judgment~\cite{tristr81}. A secondary motivation --- equally important in safety-critical deployments --- is institutional: many decisions have legal, financial, or social consequences that cannot be delegated to a machine by institutional design, regardless of the machine's empirical performance~\cite{bansal2019}.

Trist's analysis of organizational change established that technological systems and social systems are co-constituted: introducing a new technology reshapes the organizational structures around it, and those structures in turn reshape the technology's behavior and effects~\cite{tristr81}. This reciprocity applies forcefully to AI deployments. An AI system deployed without HITL mechanisms does not simply ``do the same things faster'' --- it restructures the decision-making environment in ways that can marginalize human judgment precisely when that judgment is most needed. The literature on automation bias documents this effect empirically: humans operating alongside automated systems tend to defer to automated recommendations even when those recommendations are incorrect, because the automated system's output creates a gravitational pull on the human's epistemic state~\cite{tristr81}. HITL mechanisms that position humans as passive reviewers of machine outputs are therefore not sufficient to preserve human agency in consequential decision contexts.

Effective HITL design requires that humans occupy active, structured roles that are architecturally enforced rather than socially encouraged. This means defining specific triggering conditions under which human judgment is required, providing the human with sufficient diagnostic context to exercise that judgment meaningfully, and ensuring that the human's decision propagates correctly through the system rather than being overwritten or ignored by subsequent machine processing. The canonical failure mode of HITL in deployed systems is not the absence of human involvement but the absence of the right kind of involvement --- humans presented with opaque summaries or post-hoc notifications rather than actionable escalation calls with complete diagnostic context~\cite{bansal2019}.

The Bailout Protocol operationalizes the active-HITL model by defining a specific, bounded triggering condition --- the exception state that exceeds a node's provisioned resolution capacity --- that activates the escalation pathway with full diagnostic context delivered to the human anchor. The human is not asked to review an answer; they are asked to supply one when the system has exhausted its provisioned resolution capacity. This distinction --- between HITL-as-ratification and HITL-as-escalation endpoint --- defines the protocol's contribution to the HITL design space~\cite{tristr81,bansal2019}.

% ---------------------------------------------------------------
\subsection{Fail-Safe Design in Safety-Critical Systems}
\label{sec:failsafe}

The concept of a fail-safe system --- one whose failure mode preserves human safety rather than introducing new hazards --- is well-established in safety-critical domains including aerospace, nuclear energy, and industrial process control. Classical fail-safe design proceeds by identifying the set of hazardous system states, defining the conditions under which each state is entered, and ensuring that the system transitions to a safe state whenever it cannot guarantee continued safe operation. The key property of a fail-safe system is that its safe state is not contingent on the continued correct functioning of the system itself; the fail-safe mechanism must operate even when the system is in a degraded or erroneous state~\cite{dijkstra1982,wiener1948}.

Dijkstra's formulation of structured programming and the related literature on system correctness establish a foundational principle: the reliability of a system is bounded by the reliability of its error-handling paths, not its correct-operation paths~\cite{dijkstra1982}. A system that functions correctly under nominal conditions but has unhandled or poorly-handled error modes is, in practice, less safe than a system with more limited nominal capability but robust, well-specified error behavior. This principle translates directly to autonomous AI agents: an agent whose correct-operation behavior is well-specified but whose exception-handling is underspecified represents a net safety liability that grows with the agent's operational scope.

The fail-safe principle in multi-agent systems is complicated by the absence of a single supervisory process. Unlike a safety-critical control system with a defined supervisor, a multi-agent AI network has no privileged observer who can assess global system state and invoke a global fail-safe. Wiener formalized the concept of cybernetics --- systems governed by feedback loops --- and established that control in complex systems is achieved not through central authority but through the propagation of control signals through the system's architecture~\cite{wiener1948}. In the cybernetic model, each component both receives and transmits control signals; the system is safe not because a central controller monitors everything but because every component has a defined safe-state transition that it can invoke autonomously when it detects a condition it cannot handle.

The Simplex architecture established the contemporary standard for fail-safe design in autonomous systems by introducing a dedicated safety controller that runs in parallel with a high-performance primary controller and can commandeer system execution when the primary controller's proposed action would violate a formally verified safety constraint~\cite{bansal2019}. This decoupled architecture --- performance through the primary controller, safety through an independent, verifiable safety controller --- is the direct progenitor of the BX3 three-layer model. In the Simplex model, the safety controller is a software component; in BX3, the analogous function is performed by the Fact Layer, which enforces physical and operational constraints that no Bounds Engine can override, and by the Bailout Protocol, which ensures that unresolved states reach a human safety anchor rather than being resolved by further machine inference.

What distinguishes the AI agent context from classical safety-critical control is the absence of formal safety specifications that can be mechanically verified. A safety controller in an aircraft or nuclear plant can verify proposed actions against formally specified safety invariants. An AI agent reasoning about a complex, novel directive cannot produce such a proof --- its epistemic state is probabilistic and its value function is implicit. The fail-safe mechanism in an AI agent system must therefore be calibrated not against formal safety specifications but against the boundaries of the agent's provisioned scope: the system fails safe not by guaranteeing correct behavior but by guaranteeing that any exception that exceeds its scope triggers a human-accessible escalation pathway rather than an autonomous resolution attempt~\cite{dijkstra1982,wiener1948,bansal2019}.

% ---------------------------------------------------------------
\subsection{The BX3 Framework Three-Layer Model}
\label{sec:bx3model}

The BX3 Framework defines a three-layer cognitive architecture for autonomous AI agent systems: the Purpose layer (\purpose{}), the Bounds layer (\bounds{}), and the Fact layer (\fact{}), operating over a shared Forensic Ledger maintained by the Fact layer. This architecture provides the structural foundation for the Bailout Protocol's accountability and fail-safe guarantees.

The \purpose{} layer carries intent, judgment, and final authority over any decision exceeding a node's provisioned scope. It defines the operating range $R(n)$ for each node $n$ and receives all escalated exceptions. The \purpose{} layer is responsible for mission coherence: it ensures that all actions taken by any node in the agent network remain aligned with the objectives specified by the human principal $h(n)$, a named human anchor established at node provisioning and immutable for the node's operational lifetime. The \purpose{} layer operates at the semantic level, evaluating whether proposed actions are consistent with the declared intent of the directive and flagging any action that would constitute mission drift, even if the action itself is technically within the node's operational scope. It is the first line of semantic containment.

The \bounds{} layer carries bounded reasoning and constrained execution within $R(n)$. It evaluates whether proposed actions fall within the node's authority and initiates bailout when they do not. The \bounds{} layer defines and enforces the operational envelope of each node --- the set of actions the node is permitted to take, the conditions under which it may take them, and the exceptions that must be propagated upward when conditions exceed its provisioned range. The \bounds{} layer is where the Bailout Protocol is structurally embedded: when the \bounds{} layer detects a condition that exceeds the node's provisioned resolution capacity, it triggers the protocol's escalation pathway. The \bounds{} layer is not merely a permission system; it is an active constraint-satisfaction engine that monitors operational state in real time and initiates bailout when the constraints cannot be satisfied autonomously~\cite{dijkstra1982,bansal2019}.

The \fact{} layer carries deterministic enforcement and forensic auditability. It verifies proposed actions against the authorization ledger, enforces the transition relation of the state machine, and maintains the append-only Forensic Ledger recording every protocol event as an immutable architectural record. The \fact{} layer enforces the factual grounding of all agent outputs against the Forensic Ledger --- that claims made by agents about their own state, the state of other agents, and the state of the external environment are consistent with ledger entries and are not fabricated or hallucinated. The \fact{} layer does not participate in the Bailout Protocol's decision logic; it records it. This separation of recording from decision is architecturally important: it ensures that the ledger is available as an authoritative source of truth even when the agent network is in a degraded state, including during an active bailout.

The BX3 Protocol's Immutable Parent Pointer converts the system-level accountability guarantee into a per-node architectural property. Every spawn event is treated as a delegation event with three formally defined components: a parent identity (immutable, enforced by the Fact Layer), a bounded mandate derived from the parent's Purpose Layer authorization (enforced by the Bounds Engine), and an immutable record in the Forensic Ledger (enforced by the Fact Layer). Each spawned node carries an immutable parent pointer $\alpha(n)$, and every node carries a reference to its human accountability anchor $h(n)$ --- a named human principal or institutional actor --- established at node provisioning and never modified for the node's lifetime. The Bailout Protocol exploits this structure to guarantee that exception states propagate to $h(n)$ without passing through any intermediate machine actor as a terminus.

The three-layer model maps onto the fail-safe design principles described above as follows. The \purpose{} layer supports fail-safe mechanisms by providing predefined fallback behaviors via the human anchor $h(n)$: when an action violates policy or exceeds scope, the system transitions to a safe state rather than attempting an invalid plan~\cite{wiener1948}. The \bounds{} layer supports fail-safe mechanisms by enforcing that sub-agents cannot exceed the authority granted by their parent agents, preventing cascade failures where a faulty sub-agent causes violations beyond its intended operational envelope~\cite{dijkstra1982}. The \fact{} layer supports fail-safe mechanisms by enabling precise fault localization: when an anomaly is detected, the tamper-evident identity record allows the system to pinpoint which agent produced the anomalous output and whether that agent was operating within its authorized scope~\cite{tristr81}. This alignment between the BX3 three-layer model and the fail-safe design literature establishes the conceptual substrate for the Bailout Protocol, which operationalizes these principles through specific mechanisms for detecting, attributing, and recovering from agent failures in a multi-agent environment.

The Bailout Protocol is the runtime expression of BX3's upstream accountability guarantee. Within the BX3 three-layer model, the \purpose{} layer provides the human anchor $h(n)$, the \bounds{} layer enforces the trigger condition and manages the transition to ESCALATING, and the \fact{} layer enforces the transition relation and maintains the Forensic Ledger. Together, the three layers operationalize BX3's governing principle: systems fail upward into human consciousness, never downward into algorithmic chaos.
